\section{Conclusion}
\label{Conclusion}
\setlength{\parindent}{20pt}
\justifying

\noindent
In the present study, a novel implicit bilingual statistical learning paradigm was trialled. The experimental design remains slightly too demanding for participants, but nonetheless revealed that bilingual morphemic learning can occur in the absence of explicit cues or instruction relating to the dual-language structure. A bilingual advantage was found for the detection of familiar morphemes in new input. While effect sizes were small, a possible quadrilingual disadvantage was observed, possibly due to participants processing input in a very integrated fashion instead of tracking co-occurrences within each stimulus language. The complexity of these effects, especially those arising when analysing participants past the bilingual level, demonstrates the need for refined, accurate measures of many different aspects of the multilingual experience.